\documentclass[11pt,a4paper]{article}
\usepackage[margin=27mm]{geometry}
\usepackage{amsmath,amssymb,amsthm,mathtools}
\usepackage{microtype}
\usepackage[hidelinks]{hyperref}
\newtheorem{theorem}{Theorem}
\newtheorem{lemma}[theorem]{Lemma}
\newtheorem{corollary}[theorem]{Corollary}
\newtheorem{proposition}[theorem]{Proposition}
\theoremstyle{remark}\newtheorem{remark}[theorem]{Remark}
\newcommand{\Tr}{\operatorname{Tr}}
\newcommand{\id}{\operatorname{id}}
\newcommand{\Ad}{\operatorname{Ad}}
\newcommand{\supp}{\operatorname{supp}}
\newcommand{\cQ}{\mathcal Q}
\newcommand{\cM}{\mathcal M}
\newcommand{\cN}{\mathcal N}
\newcommand{\cA}{\mathcal A}
\newcommand{\cp}{\mathrm{cp}}
\newcommand{\Herm}{\operatorname{Herm}}
\newcommand{\ket}[1]{\lvert #1\rangle}
\newcommand{\bra}[1]{\langle #1\rvert}
\newcommand{\D}{D}
\title{Nilpotent noise and complete reverse data processing}
\author{Lluis Eriksson}
\date{Working draft, 27 October 2026\\Not deposited; manuscript assessment pending}
\begin{document}
\maketitle
\begin{abstract}
We give a sufficient criterion for a reverse relative-entropy comparison
that remains valid after adjoining any finite-dimensional quantum reference.
The criterion combines a completely positive, trace-nonincreasing nilpotent
noise map, complete-positive domination on the image of a channel, and an
invertible positive output filter. A finite polynomial inverse replaces a
small-perturbation assumption and gives an explicit dimension-independent
coefficient. Applications include comparisons between qubit
amplitude-damping channels and a three-level finite-height noise family.
We also construct a rational sequence of separable two-qubit states showing
that the complete mutual-information contraction coefficient of every
nonconstant amplitude-damping channel equals one, despite its strict
ordinary relative-entropy contraction. The latter obstruction prevents an
automatic capacity application based on a strict complete contraction gap.
The reverse coefficient is not asserted sharp, and no capacity-additivity
theorem is inferred.
\end{abstract}

\section{Introduction and scope}
A reverse data-processing comparison asks how much distinguishability one
channel preserves relative to another. Ordinary comparisons optimize over
states on the input system. Retaining an arbitrary quantum reference is a
different requirement. Belzig, Gao, Smith and Wu \cite{bgsw} introduce relative
expansion coefficients, explain why complete-positive order alone does not
control the operator argument of the Bogoliubov--Kubo--Mori (BKM) metric, and
formulate complete comparisons using mutual information. Han and Li
\cite{hanli} give an exact ordinary amplitude-damping expansion coefficient;
that scalar formula is not a contribution of this paper.
Iyer, Wu, Belzig and Smith \cite{iyer} study a wider family of ordinary
Riemannian coefficients. Reverse comparisons after prior noise also relate
to the robustness questions of LaRacuente and Smith \cite{laracuente}.

Our criterion handles both the metric base point and its operator argument.
The key assumptions are an image domination
$\cQ\cM\leq_{\cp}c\cM$ and finite nilpotency $\cQ^k=0$.
The inverse of $\id+a\cQ$ is then a finite polynomial, even when
$a\sqrt c\geq1$. The conditions are finite-dimensional map conditions:
complete-positive domination can be checked on Choi operators. They are
not an assumed entropy estimate on unbounded reference dimensions.

We separately distinguish the complete mutual-information coefficient from
ordinary relative-entropy contraction and from other stabilized coefficients
\cite{hirche,bgsw}. Hirche, George, Nuradha and Wilde \cite{conditional} identify this
mutual-information coefficient with a product-reference relative-entropy
coefficient; that equivalence is prior work. It differs from coefficients
with merely equal reference marginals and from conditional coefficients
that subtract reference distinguishability.
A single-qubit reference in a separable two-qubit state suffices for the
amplitude-damping obstruction below. This does not mean that a fixed
nonzero amount of information is preserved perfectly: the information in
our witnessing sequence tends to zero.

All Hilbert spaces are finite-dimensional. All logarithms are natural.
We do not claim a solution of general nonunital-qubit additivity, a sharp
complete reverse coefficient, or a nondegradable channel with additive
quantum capacity.

\section{Metric comparisons on the positive cone}
For $S>0$ and $Y=Y^*$ define
\begin{equation}
 g_S(Y)=\int_0^\infty
 \Tr\!\left[Y(S+tI)^{-1}Y(S+tI)^{-1}\right]\,dt,
 \qquad \|Y\|_S=g_S(Y)^{1/2}.
 \label{eq:bkm}
\end{equation}
Neither $\Tr S=1$ nor $\Tr Y=0$ is required. For a semidefinite $S$ the
form is taken on its support when $\supp Y\subseteq\supp S$.
We use the associated real Hilbert norm on Hermitian operators.

\begin{lemma}[Monotonicity and order]\label{lem:metric}
Let $\Phi$ be completely positive and trace preserving. Then
$g_{\Phi(S)}(\Phi(Y))\leq g_S(Y)$ for every Hermitian $Y$.
The same inequality holds for a completely positive trace-nonincreasing
map, with the support convention above.
If $U\leq bV$, $b>0$, then
\begin{equation}
 g_U(Y)\geq b^{-1}g_V(Y).
 \label{eq:order}
\end{equation}
\end{lemma}
\begin{proof}
For a trace-one $S$, write $Y=tS+Y_0$ with $\Tr Y_0=0$.
The identities
$g_S(S,Y_0)=\Tr Y_0=0$ and $g_S(S)=1$ separate the trace component.
Differentiating relative-entropy data processing proves monotonicity on
$Y_0$; a trace-preserving map preserves the separated trace component.
Homogeneity handles arbitrary $\Tr S>0$.
For a trace-nonincreasing map append a failure flag of weight
$\Tr[(I-\Phi^*(I))S]$ to obtain a channel. The BKM form on its direct-sum
output is the desired form plus a nonnegative failure-block term.
Writing $E=I-\Phi^*(I)$, the channel is
$T\mapsto\Phi(T)\oplus\Tr(ET)$.
If $\Tr(ES)=0$, the bound $-dS\leq Y\leq dS$ forces
$\Tr(EY)=0$, so the zero failure block is omitted without a $0/0$ term.
This construction applies on the entire joint space when a reference is present.
For full-rank $S$, $-dS\leq Y\leq dS$ for some $d$, so
$\Phi(Y)$ is supported in $\Phi(S)$.

For the order statement use
$(U+tI)^{-1}\geq(bV+tI)^{-1}$ and substitute $t=bu$.
The required integrand comparison does not assume commutation:
for positive $A\geq B$,
\[
 A^{\mathsf T}\otimes A-B^{\mathsf T}\otimes B
 =(A^{\mathsf T}-B^{\mathsf T})\otimes A+
 B^{\mathsf T}\otimes(A-B)\geq0.
\]
Vectorization therefore gives
$\Tr(YAYA)\geq\Tr(YBYB)$.
For singular $U$, assume $Y$ is supported in $U$ (hence in $V$).
Compare $U+\delta I\leq bV+\delta I$, integrate, and let
$\delta\downarrow0$. The spectral formula for the BKM form gives
precisely the support-restricted limits. These are the singular cases used below.
\end{proof}

\begin{lemma}[Positive congruence]\label{lem:filter}
If $F>0$ and $F^2\geq mI$, then
\begin{equation}
 g_{FSF}(FYF)\geq m\,g_S(Y).
 \label{eq:filter}
\end{equation}
\end{lemma}
\begin{proof}
The integrand on the left of \eqref{eq:filter} is
\[
 \Tr\!\left[Y(S+tF^{-2})^{-1}Y(S+tF^{-2})^{-1}\right].
\]
Since $F^{-2}\leq m^{-1}I$, the integrand is at least its counterpart
with $S+(t/m)I$, by the same tensor-order argument as above.
Integration gives the factor $m$.
This is a comparison for an unnormalized filter; it is not an application
of channel data processing to $\Ad_F$.
\end{proof}

\section{A complete reverse comparison}
We write $\Ad_F(X)=FXF$ for $F>0$.
The order $\Phi\leq_{\cp}\Psi$ means that $\Psi-\Phi$ is completely positive.

\begin{theorem}[Nilpotent image criterion]\label{thm:main}
Let $\cM,\cN:\mathcal L(A)\to\mathcal L(B)$ be channels with $\cM(I_A)>0$.
Let $\cQ:\mathcal L(B)\to\mathcal L(B)$ be completely positive and
trace-nonincreasing, with $\cQ^k=0$ for an integer $k\geq1$.
Suppose $a\geq0$, $c>0$, $F>0$ and $m>0$ satisfy
\begin{equation}
 \cQ\cM\leq_{\cp}c\cM,\qquad
 F^2\geq mI_B,\qquad
 \cN=\Ad_F(\id+a\cQ)\cM.
 \label{eq:hyp}
\end{equation}
Set
\begin{equation}
 P_k=\sum_{j=0}^{k-1}(a\sqrt c)^j,\qquad
 \lambda=\frac{m}{(1+ac)P_k^2}.
 \label{eq:lambda}
\end{equation}
For every finite reference $R$ and density operators $\rho,\sigma$ on
$A\otimes R$ with $\supp\rho\subseteq\supp\sigma$,
\begin{equation}
 \D((\cN\otimes\id_R)(\rho)\|(\cN\otimes\id_R)(\sigma))
 \geq\lambda
 \D((\cM\otimes\id_R)(\rho)\|(\cM\otimes\id_R)(\sigma)).
 \label{eq:main}
\end{equation}
\end{theorem}
\begin{proof}
First take a full-rank $\sigma$.
Suppress tensor identities and put $S=\cM(\sigma)>0$ and $Y=\cM(X)$.
Strict positivity follows from
$\sigma\geq\epsilon I_{A\otimes R}$ and $\cM(I_A)>0$.
For any Hermitian $Z$, complete-positive domination and
Lemma~\ref{lem:metric} give
\begin{equation}
 g_S(\cQ Z)\leq c\,g_{\cQ(S)}(\cQ Z)\leq c\,g_S(Z).
 \label{eq:qnorm}
\end{equation}
Thus $\|\cQ\|_{S\to S}\leq\sqrt c$, including on nontraceless arguments.
Nilpotency supplies the exact identity
\[
 (\id+a\cQ)^{-1}=\sum_{j=0}^{k-1}(-a\cQ)^j.
\]
Writing $Z=(\id+a\cQ)Y$, the triangle inequality yields
\begin{equation}
 \|Y\|_S\leq P_k\|Z\|_S.
 \label{eq:inverse}
\end{equation}
No convergent infinite Neumann series is used.

At the base point,
\[
 \cN(\sigma)=F(S+a\cQ(S))F\leq(1+ac)FSF,
 \qquad \cN(X)=FZF.
\]
Lemmas~\ref{lem:metric} and \ref{lem:filter}, then
\eqref{eq:inverse}, imply
\begin{align}
 g_{\cN(\sigma)}(\cN(X))
 &\geq\frac{1}{1+ac}\,g_{FSF}(FZF) \notag\\
 &\geq\frac{m}{1+ac}\,g_S(Z)
 \geq\lambda g_{\cM(\sigma)}(\cM(X)).
 \label{eq:metric-main}
\end{align}
All map hypotheses and constants survive tensoring with $\id_R$.
The argument does not impose a bound on $\dim R$.

For full-rank $\rho,\sigma$, let $X=\rho-\sigma$ and
$\sigma_t=\sigma+tX$. The Hessian formula
\[
 D(\rho\|\sigma)=\int_0^1(1-t)g_{\sigma_t}(X)\,dt
\]
also holds after either channel. Integrating \eqref{eq:metric-main}
proves \eqref{eq:main}.
Finally regularize both inputs with the same full-rank state.
Output relative entropies are finite under the stated input support
inclusion. Joint convexity gives the matching upper limit for the
regularized pairs, while lower semicontinuity gives the lower limit.
This passes \eqref{eq:main} to the original pair.
\end{proof}

\begin{remark}
The domination in \eqref{eq:hyp} remains a hypothesis of this criterion.
Strict positivity of $\cM(I)$ alone is not a reverse data-processing
hypothesis. Nor is complete-positive order by itself a comparison of the
BKM operator arguments. The finite inverse in \eqref{eq:inverse} supplies
that separate step.
\end{remark}

For channels $N,M$ with the same input, define the complete
mutual-information expansion coefficient by
\begin{equation}
 \check\eta_{\mathrm{MI}}^{\mathrm{cb}}(N,M)
 =\inf_{\substack{R,\rho_{AR}\\ I(R;B_M)>0}}
       \frac{I(R;B_N)}{I(R;B_M)} .
 \label{eq:mi-coeff}
\end{equation}
The corresponding contraction coefficient uses the supremum.
These are the mutual-information quantities of \cite[Definition~7.1]{bgsw};
they are not all interchangeable with stabilized two-state coefficients.

\begin{corollary}\label{cor:mi}
Under Theorem~\ref{thm:main},
$\check\eta_{\mathrm{MI}}^{\mathrm{cb}}(\cN,\cM)\geq\lambda$.
\end{corollary}
\begin{proof}
Apply \eqref{eq:main} to $\rho_{AR}$ and
$\sigma=\rho_A\otimes\rho_R$.
The latter support contains that of $\rho_{AR}$.
Each output relative entropy is its corresponding mutual information.
\end{proof}

\section{Amplitude damping and a three-level example}
The amplitude-damping channel has Kraus operators
\[
 K_0(p)=\operatorname{diag}(1,\sqrt{1-p}),\qquad
 K_1(p)=\sqrt p\,\ket0\bra1 .
\]
For $0<p_1<p_2<1$ write $s_i=1-p_i$ and $d=p_2-p_1$.

\begin{corollary}\label{cor:ad}
For every finite reference and every supported input pair,
Theorem~\ref{thm:main} holds with
$\cM=\cA_{p_1}$, $\cN=\cA_{p_2}$ and
\begin{equation}
 \lambda_{\mathrm{AD}}
 =\frac{p_1s_2}{p_2s_1}
   \left(1+\frac{d}{\sqrt{p_1s_1}}\right)^{-2}>0.
 \label{eq:ad}
\end{equation}
The same lower bound holds in \eqref{eq:mi-coeff}.
\end{corollary}
\begin{proof}
Use
\[
 \cQ(X)=\ket0\bra1 X\ket1\bra0,\quad
 F=\operatorname{diag}(1,\sqrt{s_2/s_1}),\quad
 a=d/s_1,\quad c=s_1/p_1,\quad m=s_2/s_1.
\]
Then $\cQ^2=0$ and $\cQ$ is trace-nonincreasing.
On all matrix units,
$\cA_{p_2}=\Ad_F(\id+a\cQ)\cA_{p_1}$.
Moreover,
\[
 c\cA_{p_1}-\cQ\cA_{p_1}
   =c\,\Ad_{K_0(p_1)}
\]
is completely positive.
Substitution into \eqref{eq:lambda} gives \eqref{eq:ad}.
\end{proof}

For $p_1=1/4,p_2=1/2$ the coefficient is $1-\sqrt3/2$.
For $p_1=1/10,p_2=9/10$ it is $1/1089$, although
$a\sqrt c=8/3>1$. Neither endpoint $p_1=0$ nor $p_2=1$ is included.
Han and Li's ordinary coefficient \cite[Theorem~4]{hanli} is
$p_1s_2/(p_2s_1)$ without the extra denominator.
We make no assertion that \eqref{eq:ad} is sharp after adjoining a reference.

For a three-level example let
\[
 J=\ket0\bra1+\ket1\bra2,\qquad \cQ=\Ad_J,\qquad \cQ^3=0 .
\]
For $b>0$ set
\[
 D_b=\operatorname{diag}\bigl(1,(1+b)^{-1/2},
                                  (1+b+b^2)^{-1/2}\bigr),\qquad
 \cM_b=(\id+b\cQ+b^2\cQ^2)\Ad_{D_b}.
\]
The Kraus operators $D_b,\sqrt b JD_b,bJ^2D_b$ make $\cM_b$ a channel,
and $\cM_b(I)>0$.
The identity
\[
 b^{-1}\cM_b-\cQ\cM_b=b^{-1}\Ad_{D_b}
\]
verifies image domination.
For $0<a<1$ define
\[
 F_a=\operatorname{diag}(1,\sqrt{1-a},\sqrt{1-a+a^2}),\qquad
 \cN=\Ad_{F_a}(\id+a\cQ)\cM_b.
\]
This is a channel because
$F_a^2+aJ^*F_a^2J=I$.
Theorem~\ref{thm:main} gives the complete coefficient
\begin{equation}
 \frac{1-a}
 {(1+a/b)(1+a/\sqrt b+a^2/b)^2}.
 \label{eq:qutrit}
\end{equation}
At $b=1/100,a=1/2$, this is $1/98022$ and $a\sqrt c=5$.
This example demonstrates finite-height noise beyond the qubit case,
without claiming an optimal comparison or any capacity formula.

\section{No strict complete mutual-information contraction for amplitude damping}
The positive reverse comparison above does not imply a strict
complete contraction relative to the identity.

\begin{theorem}\label{thm:obstruction}
For every $0<p<1$,
\begin{equation}
 \eta_{\mathrm{MI}}^{\mathrm{cb}}(\cA_p,\id)=1.
 \label{eq:obstruction}
\end{equation}
A sequence of rational, full-rank separable states on two qubits attains
ratios tending to one. Its reference is not classical in any basis.
\end{theorem}
\begin{proof}
Data processing gives the upper bound one.
Let $X=\sigma_x\otimes\sigma_x$ and, for integers $n\geq3$, put
\[
 \begin{gathered}
 x_n=1/n,\qquad r_n=2^{-n^2},\qquad \epsilon_n=r_n/n^2,\\
 S_n=\operatorname{diag}(1-x_n,x_n)
        \otimes\operatorname{diag}(1-r_n,r_n),\qquad
 \rho_n=S_n+\epsilon_n X.
\end{gathered}
\]
All entries are rational, both marginal states are unchanged, and
$\lambda_{\min}(S_n)=r_n/n$.
Since $\|X\|_\infty=1$,
\begin{equation}
 (1-1/n)S_n\leq\rho_n\leq(1+1/n)S_n.
 \label{eq:finite-order}
\end{equation}
In particular $\rho_n$ is a full-rank state.

For positive numbers write
$\ell(a,b)=(\log a-\log b)/(a-b)$, with $\ell(a,a)=1/a$.
At the product base
$S_{x,r}=\operatorname{diag}(1-x,x)\otimes\operatorname{diag}(1-r,r)$,
let $s=1-p$, $u=1-sx$ and $v=sx$.
Direct evaluation of \eqref{eq:bkm} gives
\begin{align}
 g_{\mathrm{in}}
 &=2\{\ell((1-x)(1-r),xr)+\ell((1-x)r,x(1-r))\},\label{eq:gin}\\
 g_{\mathrm{out}}
 &=2s\{\ell(u(1-r),vr)+\ell(ur,v(1-r))\}.
 \label{eq:gout}
\end{align}
The factor $s$ comes from $\cA_p(\sigma_x)=\sqrt s\,\sigma_x$.

For fixed $x$, letting $r\downarrow0$ in these forms gives the ratio
$(1-x)/(1-sx)$.
For the particular joint sequence $x_n,r_n$, the same limit to one is
seen directly. Put $T_n=n^2\log2$ and, for fixed $d>0$, $q=d/n$.
The sum of the two log means at base populations $(1-q,q)$ is
\[
 H_n(q)=\frac{T_n+\log((1-q)/q)+\log(1-r_n)}{1-q-r_n}
       +\frac{T_n-\log((1-q)/q)+\log(1-r_n)}{q-r_n}.
\]
Since $\log((1-q)/q)=\log n-\log d+O_d(n^{-1})$,
the first term is $O(T_n)$ and the second is
$(n/d)T_n(1+o(1))$. Taking $d=1$ and $d=s$ shows
\[
 g_{\mathrm{in}}(x_n,r_n)\sim2n^3\log2,\qquad
 g_{\mathrm{out}}(x_n,r_n)\sim2n^3\log2
\]
for each fixed $0<s<1$.

It remains to control the finite mutual informations, rather than merely
their Hessians.
The bounds \eqref{eq:finite-order} hold along the segment from $S_n$ to
$\rho_n$ and after applying the channel.
Metric order and the Hessian integral imply, at input and at output,
\[
 \frac{\epsilon_n^2}{2(1+1/n)}g
 \leq D(\rho_n\|S_n)
 \leq\frac{\epsilon_n^2}{2(1-1/n)}g,
\]
using the appropriate output pair and form in the output expression.
These relative entropies are the mutual informations because the marginal
states are unchanged. Their ratio is at least
\[
 \frac{1-1/n}{1+1/n}
 \frac{g_{\mathrm{out}}(x_n,r_n)}{g_{\mathrm{in}}(x_n,r_n)}
 \longrightarrow1.
\]
This proves \eqref{eq:obstruction} without interchanging an uncontrolled
Taylor approximation and a boundary limit.

To prove separability, abbreviate $x=x_n,r=r_n,\epsilon=\epsilon_n$ and set
\[
 D_A=\operatorname{diag}(\sqrt{1-x},\sqrt x),\qquad
 D_R=\operatorname{diag}(\sqrt{1-r},\sqrt r),\qquad
 t=\frac{\epsilon}{\sqrt{x(1-x)r(1-r)}}\in(0,1).
\]
Then
\[
 \rho_n=(D_A\otimes D_R)(I_4+t\,\sigma_x\otimes\sigma_x)(D_A\otimes D_R).
\]
If $P_a$ denotes the $\sigma_x$ eigenprojector with eigenvalue $a\in\{-1,1\}$,
the middle operator equals
$\sum_{a,b}(1+tab)P_a\otimes P_b$.
All coefficients are positive. Local filtering preserves the separable
positive cone, giving a separable decomposition of the normalized state.
Finally the reference marginal has distinct eigenvalues and
$[\rho_n,I\otimes(\rho_n)_R]\neq0$.
Every state classical on $R$ commutes with its reference marginal, so
$\rho_n$ is not classical on $R$ in any basis.
\end{proof}

\begin{remark}[Small information, not perfect transmission]
The input mutual information in this sequence tends to zero; the above
bounds give
$I(A;R)_{\rho_n}\sim(\log2)\,2^{-2n^2}/n$.
The theorem concerns a supremum of ratios. It is not a claim of noiseless
transmission of a fixed amount of information.
At $p=1$ the output is constant and the mutual-information coefficient is
zero, so that endpoint is not part of Theorem~\ref{thm:obstruction}.
\end{remark}

Ordinary contraction is different: the bound
$\eta(\cA_p)\leq\sqrt{1-p}<1$ is given in
\cite[Lemma~6.3]{bgsw}.
For a classical auxiliary the mutual information is an average of relative
entropies, so the same strict ordinary coefficient controls its contraction.
Theorem~\ref{thm:obstruction} shows that separable nonclassical references
can invalidate that strict coefficient.

For a degradable amplitude-damping main channel with damping
$\gamma\in(0,1/2)$, its degrading map has damping
$(1-2\gamma)/(1-\gamma)\in(0,1)$.
Its complete mutual-information contraction coefficient is therefore one.
Consequently the factor $1-\eta_{\mathrm{MI}}^{\mathrm{cb}}$ in the
sufficient construction discussed in \cite[Section~7]{bgsw} vanishes.
Our positive relative-expansion bound alone does not establish the
informationally degradable nondegradable example sought by that criterion.

\section{Verification, limits and review context}
The proofs use finite map identities, BKM monotonicity, operator order and
relative-entropy data processing. Algebraic checks verify the amplitude
and three-level map identities, positivity of selected rational witnesses,
and the stated example constants. Numerical checks at selected correlated
base states examine whole finite Hermitian tangent spaces, including
nontraceless arguments. Neither kind of finite check is a replacement for
the universal arguments or a certificate of worldwide novelty.

The standard ordinary coefficient formulas are attributed to their sources.
The image domination, nilpotency, full-output-support and positive-filter
hypotheses of Theorem~\ref{thm:main} remain restrictions.
The result does not address all pairs of channels, arbitrary
infinite-dimensional references, or a sharp value of the complete coefficient.
It does not solve the separate non-GAD output-entropy target considered
in the author's earlier working notes.

\subsection*{Model assessments and review context}
A separate mathematical component audit (audit917, 27 October 2026) was
requested with configuration GPT-6 Astra, reasoning effort xhigh, on a
hash-bound version of the underlying argument and checker. Its terminal
report found the structural theorem, both channel examples and the
separable mutual-information obstruction valid under the stated
hypotheses, with no fatal proof error found. It assigned no numerical grade.
The complete report is retained with SHA-256
\texttt{e5ca2f4dc39ca5de352d72e7751bafdb159a2}\allowbreak
\texttt{acf359121fb40f650c743444259}.
The audit required broader prior-art attribution, explicit singular-support
and failure-flag details, and careful separation of coefficient definitions.
Those points are incorporated here. Its reservations remain: novelty is
unresolved; ordinary amplitude-damping positivity and exact coefficients
are prior work; the complete lower bound is not proved optimal; no
capacity or additivity theorem follows. A local algebraic issue reported
in the rendered proof of an antecedent is not used to discredit that
antecedent's proposition or establish priority.
The requested model identity and effort are not externally attested;
the context was separate but the filesystem was shared. Independence and
runtime isolation are unverified.
A component audit is not an assessment of this exact PDF.
Any subsequent manuscript evaluation must be read with its exact-PDF hash,
full report, objections and context; no numerical grade is asserted here.

\subsection*{AI and responsibility disclosure}
The named author directed substantial OpenAI Codex assistance in derivation,
writing and numerical work. The requested component audit used a separate
conversational context with a shared filesystem. Effective model identity,
effort, runtime isolation and independence have not been independently
attested. No independent human review, Lean formalization, certified
worldwide originality or predicted assessment score is claimed.

\begin{thebibliography}{9}
\bibitem{bgsw}
P. Belzig, L. Gao, G. Smith and P. Wu,
\emph{Reverse-type Data Processing Inequality},
arXiv:2411.19890v3 (2025).
\url{https://arxiv.org/abs/2411.19890v3}.

\bibitem{hanli}
J. Han and Y. Li,
\emph{Convergence Analysis of the Quantum Blahut--Arimoto Algorithm and the
Contraction Coefficient of Qubit Amplitude-Damping Channels},
arXiv:2610.29373v1 (2026).
\url{https://arxiv.org/abs/2610.29373v1}.

\bibitem{iyer}
S. Iyer, P. Wu, P. Belzig and G. Smith,
\emph{Quantum $f$-divergences and Their Local Behaviour:
An Analysis via Relative Expansion Coefficients},
arXiv:2510.06183v2 (2025).
\url{https://arxiv.org/abs/2510.06183v2}.

\bibitem{hirche}
C. Hirche, C. Rouz\'e and D. Stilck Fran\c{c}a,
\emph{On contraction coefficients, partial orders and approximation of
capacities for quantum channels}, Quantum \textbf{6}, 862 (2022).
\url{https://doi.org/10.22331/q-2022-11-28-862}.

\bibitem{conditional}
C. Hirche, I. George, T. Nuradha and M. M. Wilde,
\emph{Conditional contraction coefficients and their applications to quantum networks},
arXiv:2608.27171v1 (2026).
\url{https://arxiv.org/abs/2608.27171v1}.

\bibitem{laracuente}
N. LaRacuente and G. Smith,
\emph{Information Fragility or Robustness Under Quantum Channels},
arXiv:2312.17450v1 (2023).
\url{https://arxiv.org/abs/2312.17450v1}.
\end{thebibliography}
\end{document}
